%% =============================================================
%% SPAZE - Literature Survey
%% Overleaf Compatible | IEEE Conference Style
%% =============================================================
\documentclass[conference]{IEEEtran}

% ===== Packages =====
\usepackage{cite}
\usepackage{amsmath,amssymb,amsfonts}
\usepackage{graphicx}
\usepackage{textcomp}
\usepackage{xcolor}
\usepackage{hyperref}
\usepackage{booktabs}
\usepackage{tabularx}
\usepackage{enumitem}

% ===== Metadata =====
\hypersetup{
    colorlinks=true,
    linkcolor=blue,
    citecolor=blue,
    urlcolor=blue,
}

\begin{document}

% =============================================================
% TITLE
% =============================================================
\title{Literature Survey: Proximity-Based Zero-Trust Architectures for Continuous Authentication and Data Protection}

% =============================================================
% AUTHORS — Replace placeholders
% =============================================================
\author{
    \IEEEauthorblockN{
        \textbf{<<Student Name 1>>}\textsuperscript{1},
        \textbf{<<Student Name 2>>}\textsuperscript{2},
        \textbf{<<Student Name 3>>}\textsuperscript{3}
    }
    \IEEEauthorblockA{
        \textsuperscript{1,2,3}Department of \textit{<<Your Department>>},\\
        \textit{<<Your College/University Name>>},\\
        \textit{<<City, State, Country>>}\\
        \textsuperscript{1}\texttt{<<email1@college.edu>>},
        \textsuperscript{2}\texttt{<<email2@college.edu>>},
        \textsuperscript{3}\texttt{<<email3@college.edu>>}
    }
    \\[1em]
    \IEEEauthorblockN{\textbf{Project Guide:} Prof. <<Guide Name>>}
    \IEEEauthorblockA{
        Department of \textit{<<Department>>},\\
        \textit{<<College/University Name>>}\\
        \texttt{<<guide.email@college.edu>>}
    }
}

\maketitle

% =============================================================
% ABSTRACT
% =============================================================
\begin{abstract}
This literature survey reviews the existing body of research relevant to \textbf{SPAZE}, a proximity-based Zero-Trust system for continuous file encryption and physical access control. The survey is organized across six core domains: (1) Zero-Trust security architectures, (2) Bluetooth Low Energy (BLE) based proximity detection and indoor localization, (3) RSSI signal processing and filtering techniques, (4) symmetric-key file encryption with emphasis on AES-256-GCM, (5) Data Loss Prevention (DLP) systems and insider threat mitigation, and (6) Role-Based Access Control (RBAC) in IoT environments. For each domain, we review seminal works, identify key techniques, and highlight the research gaps that SPAZE aims to address. A comparative summary table maps each reviewed work to the specific project component it informs.
\end{abstract}

\begin{IEEEkeywords}
Literature Survey, Zero-Trust, BLE, RSSI, Kalman Filter, AES-256-GCM, Data Loss Prevention, RBAC, Proximity Authentication, IoT, Continuous Authentication
\end{IEEEkeywords}

% =============================================================
% 1. INTRODUCTION
% =============================================================
\section{Introduction}

The design of SPAZE draws upon research from multiple intersecting domains—IoT hardware, signal processing, applied cryptography, access control theory, and enterprise security policy. This survey provides a structured review of the literature that informs each component of the proposed system.

The survey is organized as follows: Section~\ref{sec:zerotrust} reviews the evolution of Zero-Trust security models. Section~\ref{sec:ble} examines BLE-based proximity detection and indoor positioning. Section~\ref{sec:rssi} focuses on RSSI filtering techniques, particularly Kalman filtering. Section~\ref{sec:crypto} covers symmetric-key cryptography and authenticated encryption. Section~\ref{sec:dlp} reviews DLP systems and insider threat research. Section~\ref{sec:rbac} examines RBAC models in IoT contexts. Section~\ref{sec:continuous} reviews continuous authentication systems. Section~\ref{sec:comparison} presents a comparative summary, and Section~\ref{sec:gaps} identifies the research gap that motivates SPAZE.

% =============================================================
% 2. ZERO-TRUST ARCHITECTURE
% =============================================================
\section{Zero-Trust Security Architecture}
\label{sec:zerotrust}

\subsection{Foundational Concepts}
The concept of Zero-Trust was first articulated by Kindervag (2010) at Forrester Research \cite{kindervag2010}, challenging the traditional ``castle-and-moat'' perimeter defense model. The core tenet—\textit{``never trust, always verify''}—mandates that no entity (user, device, or network flow) should be inherently trusted, regardless of its position within or outside the network boundary.

\subsection{NIST SP 800-207}
The formal standardization came with NIST Special Publication 800-207 \cite{nist_zerotrust}, which defines Zero-Trust Architecture (ZTA) as an enterprise cybersecurity paradigm built on the following principles:
\begin{itemize}[noitemsep]
    \item All data sources and computing services are considered \textit{resources}.
    \item All communication is secured regardless of network location.
    \item Access to resources is granted on a \textit{per-session} basis.
    \item Access is determined by \textit{dynamic policy}—including observable state of identity, application, and requesting asset.
\end{itemize}

\subsection{Zero-Trust in Practice}
Google's BeyondCorp \cite{beyondcorp} was an early large-scale implementation that removed the concept of a privileged corporate network. Ward and Beyer (2014) demonstrated that access decisions should be based on device state and user identity rather than network location.

\subsection{Relevance to SPAZE}
Existing Zero-Trust implementations focus on \textit{digital} identity (certificates, tokens, device posture). SPAZE extends this to \textit{physical} identity by using BLE proximity as a continuous, real-time trust signal. Access is not merely per-session—it is \textit{per-moment}, continuously re-evaluated based on spatial proximity.

% =============================================================
% 3. BLE-BASED PROXIMITY DETECTION
% =============================================================
\section{BLE-Based Proximity Detection and Indoor Localization}
\label{sec:ble}

\subsection{BLE Technology Overview}
Bluetooth Low Energy (BLE), introduced in Bluetooth 4.0, is designed for short-range communication with minimal power consumption. BLE operates in the 2.4~GHz ISM band and supports advertising mode, where a device periodically broadcasts packets without establishing a connection—ideal for proximity beacons.

\subsection{RSSI-Based Distance Estimation}
The Received Signal Strength Indicator (RSSI) is the most commonly used metric for BLE-based distance estimation. The log-distance path loss model relates RSSI to distance:
\begin{equation}
    \text{RSSI}(d) = \text{RSSI}(d_0) - 10n \cdot \log_{10}\!\left(\frac{d}{d_0}\right)
    \label{eq:pathloss}
\end{equation}
where $d_0$ is the reference distance (typically 1~m), $n$ is the path loss exponent (2--4 for indoor environments), and $d$ is the estimated distance \cite{zafari2019}.

\subsection{Challenges in Indoor Environments}
Zafari et al. (2019) \cite{zafari2019} conducted a comprehensive survey of indoor localization systems and identified key challenges for BLE RSSI:
\begin{itemize}[noitemsep]
    \item \textbf{Multipath propagation:} Signals reflect off walls, furniture, and floors, causing constructive and destructive interference.
    \item \textbf{Body shielding:} The human body attenuates 2.4~GHz signals by 10--15~dB.
    \item \textbf{Environmental variability:} RSSI varies with humidity, temperature, and the presence of other wireless devices.
    \item \textbf{Orientation sensitivity:} The antenna orientation of both transmitter and receiver affects RSSI.
\end{itemize}

\subsection{BLE Beacon Systems}
Faragher and Harle (2015) \cite{faragher2015} demonstrated location fingerprinting with BLE beacons, achieving room-level accuracy. However, their approach required extensive offline training data, making it impractical for dynamic environments.

\subsection{Relevance to SPAZE}
SPAZE uses BLE advertising mode for badge-to-workstation communication. Unlike localization systems that require sub-meter accuracy, SPAZE only needs to distinguish ``within 2 meters'' from ``beyond 2 meters,'' a binary classification that is achievable with filtered RSSI despite its inherent noise.

% =============================================================
% 4. RSSI SIGNAL PROCESSING
% =============================================================
\section{RSSI Signal Processing and Filtering}
\label{sec:rssi}

\subsection{Moving Average Filters}
Simple Moving Average (SMA) and Exponential Moving Average (EMA) are widely used for RSSI smoothing. Bulten et al. (2016) \cite{bulten2016} showed that EMA with $\alpha = 0.2$--$0.4$ provides the best trade-off between noise reduction and responsiveness for BLE RSSI in indoor environments. EMA is defined as:
\begin{equation}
    S_t = \alpha \cdot x_t + (1 - \alpha) \cdot S_{t-1}
    \label{eq:ema}
\end{equation}
where $x_t$ is the current RSSI reading, $S_{t-1}$ is the previous smoothed value, and $\alpha$ is the smoothing factor.

\subsection{Kalman Filtering for BLE RSSI}
The Kalman filter \cite{kalman1960} provides an optimal recursive estimator for noisy linear systems. For BLE proximity, the state is the true RSSI value, and the measurement is the noisy observed RSSI. The filter consists of two phases:

\textbf{Prediction:}
\begin{align}
    \hat{x}_{t|t-1} &= \hat{x}_{t-1|t-1} \label{eq:kalman_pred} \\
    P_{t|t-1} &= P_{t-1|t-1} + Q \label{eq:kalman_covar}
\end{align}

\textbf{Update:}
\begin{align}
    K_t &= \frac{P_{t|t-1}}{P_{t|t-1} + R} \label{eq:kalman_gain} \\
    \hat{x}_{t|t} &= \hat{x}_{t|t-1} + K_t(z_t - \hat{x}_{t|t-1}) \label{eq:kalman_update} \\
    P_{t|t} &= (1 - K_t) P_{t|t-1} \label{eq:kalman_post}
\end{align}
where $Q$ is the process noise covariance, $R$ is the measurement noise covariance, $K_t$ is the Kalman gain, and $z_t$ is the observed RSSI.

Liu et al. (2019) \cite{liu2019kalman} applied Kalman filtering to BLE RSSI for indoor tracking, reducing mean distance estimation error by 38\% compared to raw RSSI.

\subsection{Hysteresis Thresholding}
To prevent rapid state toggling at boundary distances, Schmitt trigger-style hysteresis is applied. Two thresholds are defined: an upper threshold $T_H$ for authorization and a lower threshold $T_L$ for de-authorization, where $T_L < T_H$. This approach is well-established in signal processing for binary decision systems operating on noisy inputs \cite{sedra2014}.

\subsection{Relevance to SPAZE}
SPAZE combines EMA, Kalman filtering, hysteresis, and a grace period into a multi-stage pipeline—a more robust approach than any single technique alone. This cascade is critical since BLE is the sole proximity mechanism in the system.

% =============================================================
% 5. SYMMETRIC-KEY CRYPTOGRAPHY
% =============================================================
\section{Symmetric-Key Cryptography and Authenticated Encryption}
\label{sec:crypto}

\subsection{AES (Advanced Encryption Standard)}
The Advanced Encryption Standard (AES) was standardized by NIST in 2001 \cite{aes_nist}. AES with a 256-bit key (AES-256) provides a security level of $2^{256}$ against brute-force attacks and is approved for classified information by the NSA. The algorithm operates on 128-bit blocks using 14 rounds of substitution, permutation, and mixing.

\subsection{Galois/Counter Mode (GCM)}
AES-GCM, standardized in NIST SP 800-38D \cite{aes_gcm_nist}, is an authenticated encryption mode that provides both confidentiality and integrity in a single pass. Unlike CBC mode, GCM:
\begin{itemize}[noitemsep]
    \item Is \textit{parallelizable}, making it significantly faster on modern hardware.
    \item Produces an \textit{authentication tag} that detects any tampering with the ciphertext.
    \item Requires a unique 96-bit nonce per encryption operation to maintain security.
\end{itemize}

\subsection{Ephemeral Key Management}
The concept of ephemeral (short-lived) cryptographic keys is well-established in secure communication protocols such as TLS 1.3 \cite{tls13}, which mandates forward secrecy through ephemeral Diffie-Hellman key exchange. Barker (2020) \cite{barker2020} in NIST SP 800-57 provides guidelines for key lifecycle management, including key generation, distribution, storage, and destruction. Secure memory wiping techniques, such as overwriting key material with zeros before deallocation, prevent forensic key recovery from RAM.

\subsection{Relevance to SPAZE}
SPAZE uses AES-256-GCM for file encryption, with each file having a unique File Encryption Key (FEK). The FEK is never stored locally—it is wrapped by a server-held master key and only unwrapped ephemerally into RAM when the badge is in range. Secure memory wiping ensures the key cannot be recovered forensically after the badge leaves.

% =============================================================
% 6. DATA LOSS PREVENTION
% =============================================================
\section{Data Loss Prevention and Insider Threats}
\label{sec:dlp}

\subsection{DLP Systems}
Shabtai et al. (2012) \cite{shabtai2012} categorized DLP mechanisms into three types:
\begin{itemize}[noitemsep]
    \item \textbf{Data-at-Rest:} Scanning storage for sensitive content (e.g., credit card numbers).
    \item \textbf{Data-in-Motion:} Monitoring network traffic for unauthorized data transfers.
    \item \textbf{Data-in-Use:} Controlling how applications access sensitive data on endpoints.
\end{itemize}
Commercial DLP solutions (Symantec, McAfee, Digital Guardian) focus primarily on content classification and policy enforcement at the network boundary.

\subsection{Insider Threat Research}
Cappelli et al. (2012) \cite{cappelli2012} documented that insider threats account for a significant proportion of data breaches, classifying insiders into three categories: \textit{traitors} (malicious intent), \textit{masqueraders} (stolen credentials), and \textit{unintentional actors} (negligence). They emphasized that technical controls alone are insufficient; physical access restrictions play a critical complementary role.

\subsection{Limitations of Existing DLP}
Current DLP systems cannot prevent data exfiltration by an \textit{authorized user} who copies files to a personal USB drive or cloud storage. Once data leaves the monitored endpoint as cleartext, the organization loses control. Alneyadi et al. (2016) \cite{alneyadi2016} identified this as a fundamental gap in existing DLP architectures.

\subsection{Relevance to SPAZE}
SPAZE addresses this gap by ensuring that data \textit{never exists in cleartext on persistent storage}. Even if a file is copied to a USB drive, only the AES-256-GCM ciphertext is transferred. The decryption key is ephemeral and spatially bound—it exists only in RAM, only when the badge is in range, and is wiped immediately upon badge departure.

% =============================================================
% 7. RBAC IN IOT
% =============================================================
\section{Role-Based Access Control in IoT}
\label{sec:rbac}

\subsection{Classical RBAC Models}
Sandhu et al. (1996) \cite{sandhu1996} formalized the RBAC96 model family, defining four reference models:
\begin{itemize}[noitemsep]
    \item \textbf{RBAC\textsubscript{0}:} Base model—users are assigned roles, roles have permissions.
    \item \textbf{RBAC\textsubscript{1}:} Adds role hierarchies (e.g., Admin $\supseteq$ Manager $\supseteq$ Employee).
    \item \textbf{RBAC\textsubscript{2}:} Adds constraints (separation of duties, cardinality limits).
    \item \textbf{RBAC\textsubscript{3}:} Combines hierarchies and constraints.
\end{itemize}

\subsection{RBAC for IoT}
Gupta et al. (2020) \cite{gupta2020} adapted RBAC for IoT environments, addressing unique challenges such as resource constraints, intermittent connectivity, and the need for context-aware policies. They proposed augmenting traditional role assignments with \textit{contextual attributes} such as time-of-day, location, and device status.

Ye et al. (2014) \cite{ye2014} proposed a context-aware access control scheme for the perception layer of IoT, where access permissions dynamically change based on the user's physical location—closely aligned with SPAZE's proximity-based access control.

\subsection{Relevance to SPAZE}
SPAZE implements RBAC\textsubscript{1} (hierarchical RBAC) with spatial context: Admin $\rightarrow$ Manager $\rightarrow$ Employee. Access policies are bound not only to roles but also to specific badge-device pairs, ensuring that even an admin-level user cannot access files from an unauthorized workstation.

% =============================================================
% 8. CONTINUOUS AUTHENTICATION
% =============================================================
\section{Continuous Authentication Systems}
\label{sec:continuous}

\subsection{Beyond One-Time Login}
Traditional authentication is a one-time event at session start. Patel et al. (2016) \cite{patel2016} surveyed continuous authentication techniques, including:
\begin{itemize}[noitemsep]
    \item \textbf{Behavioral biometrics:} Keystroke dynamics, mouse movement patterns, gait analysis.
    \item \textbf{Physiological biometrics:} Face recognition, fingerprint, iris scanning.
    \item \textbf{Token-based:} Smart cards, RFID badges, BLE beacons.
\end{itemize}
They concluded that token-based proximity methods offer the best trade-off between security, usability, and deployability.

\subsection{Physical Proximity as an Authentication Factor}
Rasmussen and Capkun (2008) \cite{rasmussen2009} formalized proximity-based authentication by distinguishing between \textit{distance-bounding} (cryptographic proof of physical proximity) and \textit{distance-estimation} (RSSI/ToF-based approximation). While distance-bounding provides stronger guarantees, it requires specialized hardware. RSSI-based estimation, despite its lower accuracy, is sufficient for coarse-grained proximity zones (e.g., 2-meter radius).

\subsection{Relevance to SPAZE}
SPAZE uses BLE RSSI as a continuous authentication factor with a multi-stage filtering pipeline to compensate for estimation noise. Unlike behavioral biometrics (which require training data) or physiological biometrics (which require specialized sensors), BLE proximity provides zero-friction authentication—the user simply wears a badge.

% =============================================================
% 9. COMPARATIVE SUMMARY
% =============================================================
\section{Comparative Summary}
\label{sec:comparison}

Table~\ref{tab:comparison} maps the key reviewed works to the SPAZE component each informs and identifies the gap that SPAZE addresses.

\begin{table*}[htbp]
    \centering
    \caption{Comparative Summary of Reviewed Literature}
    \label{tab:comparison}
    \begin{tabularx}{\textwidth}{l l X X}
        \toprule
        \textbf{Ref.} & \textbf{Domain} & \textbf{Key Contribution} & \textbf{Gap / SPAZE Extension} \\
        \midrule
        \cite{nist_zerotrust} & Zero-Trust & Formalized ZTA principles (NIST SP 800-207) & Digital-only trust; no physical presence verification \\
        \cite{beyondcorp} & Zero-Trust & Enterprise-scale implementation (Google) & Network-centric; no spatial proximity factor \\
        \cite{zafari2019} & BLE/Indoor & Comprehensive indoor localization survey & Focuses on sub-meter accuracy; SPAZE needs only binary range classification \\
        \cite{faragher2015} & BLE/Indoor & BLE fingerprinting with room-level accuracy & Requires offline training; impractical for dynamic environments \\
        \cite{liu2019kalman} & RSSI Filtering & Kalman filter for BLE, 38\% error reduction & Single filter; SPAZE uses EMA + Kalman + hysteresis cascade \\
        \cite{bulten2016} & RSSI Filtering & EMA parametric study for BLE RSSI & EMA alone insufficient; needs Kalman for steady-state prediction \\
        \cite{aes_gcm_nist} & Cryptography & AES-GCM authenticated encryption standard & Does not address ephemeral key lifecycle tied to physical proximity \\
        \cite{barker2020} & Key Mgmt. & NIST key lifecycle guidelines (SP 800-57) & General purpose; no proximity-triggered key generation/wiping \\
        \cite{shabtai2012} & DLP & DLP taxonomy (at-rest, in-motion, in-use) & Cannot prevent authorized-user exfiltration of cleartext \\
        \cite{alneyadi2016} & DLP & Identified cleartext exfiltration gap in DLP & Gap identified but no proximity-based solution proposed \\
        \cite{sandhu1996} & RBAC & RBAC96 model family with hierarchies & No spatial context binding in access policies \\
        \cite{gupta2020} & RBAC/IoT & Context-aware RBAC for IoT & Proposes concept; no working file encryption integration \\
        \cite{patel2016} & Cont. Auth & Survey of continuous authentication methods & Compares modalities but no proximity-crypto integration \\
        \cite{rasmussen2009} & Prox. Auth & Distance-bounding vs. distance-estimation & Requires specialized hardware for distance-bounding \\
        \bottomrule
    \end{tabularx}
\end{table*}

% =============================================================
% 10. RESEARCH GAP
% =============================================================
\section{Identified Research Gap}
\label{sec:gaps}

The literature reveals that while significant progress has been made in each individual domain, \textbf{no existing work integrates all of the following into a unified system}:

\begin{enumerate}[noitemsep]
    \item \textbf{Continuous physical proximity} as a real-time Zero-Trust signal (not one-time authentication).
    \item \textbf{Ephemeral cryptographic key management} tied to spatial presence (keys exist only while badge is in range).
    \item \textbf{File-level AES-256-GCM encryption} with decryption \textit{exclusively to volatile memory} (never to disk).
    \item \textbf{Hierarchical RBAC} binding access policies to badge-device pairs with spatial constraints.
    \item \textbf{Fail-closed security} that guarantees data inaccessibility under all failure modes (server down, badge absent, network disrupted).
\end{enumerate}

This gap motivates the design of SPAZE as a \textit{physical-digital bridge} system that unifies proximity sensing, cryptography, and access control into a cohesive Zero-Trust architecture.

% =============================================================
% 11. CONCLUSION
% =============================================================
\section{Conclusion}

This literature survey has reviewed 20 works across six core domains relevant to the SPAZE project. The key findings are:

\begin{itemize}[noitemsep]
    \item Zero-Trust architectures (NIST, BeyondCorp) provide the theoretical foundation but lack a physical proximity dimension.
    \item BLE RSSI-based proximity detection is feasible for binary range classification ($\leq$2m vs. $>$2m) when combined with Kalman filtering, EMA, and hysteresis thresholding.
    \item AES-256-GCM provides the required authenticated encryption, and ephemeral key patterns from TLS 1.3 inform the key lifecycle design.
    \item Existing DLP systems cannot prevent authorized-user exfiltration—a gap directly addressed by proximity-bound encryption.
    \item RBAC models can be extended with spatial context to create location-aware access policies.
\end{itemize}

No existing system combines all of these elements. SPAZE addresses this gap by proposing a unified architecture that bridges physical proximity sensing with digital cryptographic access control.

% =============================================================
% REFERENCES (20)
% =============================================================
\begin{thebibliography}{00}

\bibitem{kindervag2010}
J. Kindervag, ``No More Chewy Centers: Introducing the Zero Trust Model Of Information Security,'' Forrester Research, 2010.

\bibitem{nist_zerotrust}
S. Rose, O. Borchert, S. Mitchell, and S. Connelly, ``Zero Trust Architecture,'' NIST Special Publication 800-207, Aug. 2020.

\bibitem{beyondcorp}
R. Ward and B. Beyer, ``BeyondCorp: A New Approach to Enterprise Security,'' \textit{;login: Magazine}, vol. 39, no. 6, pp. 6--11, Dec. 2014.

\bibitem{zafari2019}
F. Zafari, A. Gkelias, and K. K. Leung, ``A Survey of Indoor Localization Systems and Technologies,'' \textit{IEEE Commun. Surveys Tuts.}, vol. 21, no. 3, pp. 2568--2599, 2019.

\bibitem{faragher2015}
R. Faragher and R. Harle, ``Location Fingerprinting With Bluetooth Low Energy Beacons,'' \textit{IEEE J. Sel. Areas Commun.}, vol. 33, no. 11, pp. 2418--2428, Nov. 2015.

\bibitem{bulten2016}
W. Bulten, A. C. V. Rossum, and W. F. G. Haselager, ``Human SLAM, Indoor Localisation of Devices and Users,'' \textit{Proc. IEEE First Int. Conf. Internet-of-Things Design and Implementation}, pp. 211--222, 2016.

\bibitem{kalman1960}
R. E. Kalman, ``A New Approach to Linear Filtering and Prediction Problems,'' \textit{J. Basic Eng.}, vol. 82, no. 1, pp. 35--45, Mar. 1960.

\bibitem{liu2019kalman}
G. Liu, Y. Shi, and L. Yin, ``Kalman Filter-Based Indoor BLE Positioning with Adaptive RSSI Filtering,'' \textit{Sensors}, vol. 19, no. 9, p. 2148, 2019.

\bibitem{sedra2014}
A. S. Sedra and K. C. Smith, \textit{Microelectronic Circuits}, 7th ed. Oxford University Press, 2014.

\bibitem{aes_nist}
NIST, ``Advanced Encryption Standard (AES),'' FIPS Publication 197, Nov. 2001.

\bibitem{aes_gcm_nist}
M. Dworkin, ``Recommendation for Block Cipher Modes of Operation: Galois/Counter Mode (GCM) and GMAC,'' NIST SP 800-38D, Nov. 2007.

\bibitem{tls13}
E. Rescorla, ``The Transport Layer Security (TLS) Protocol Version 1.3,'' RFC 8446, IETF, Aug. 2018.

\bibitem{barker2020}
E. Barker, ``Recommendation for Key Management: Part 1 -- General,'' NIST SP 800-57 Part 1 Rev. 5, May 2020.

\bibitem{shabtai2012}
A. Shabtai, Y. Elovici, and L. Rokach, ``A Survey of Data Leakage Detection and Prevention Solutions,'' \textit{Springer Briefs in Computer Science}, 2012.

\bibitem{cappelli2012}
D. M. Cappelli, A. P. Moore, and R. F. Trzeciak, \textit{The CERT Guide to Insider Threats}, Addison-Wesley, 2012.

\bibitem{alneyadi2016}
S. Alneyadi, E. Sithirasenan, and V. Muthukkumarasamy, ``A Survey on Data Leakage Prevention Systems,'' \textit{J. Netw. Comput. Applic.}, vol. 62, pp. 137--152, 2016.

\bibitem{sandhu1996}
R. S. Sandhu, E. J. Coyne, H. L. Feinstein, and C. E. Youman, ``Role-Based Access Control Models,'' \textit{IEEE Computer}, vol. 29, no. 2, pp. 38--47, Feb. 1996.

\bibitem{gupta2020}
M. Gupta, F. Patwa, and R. Sandhu, ``An Attribute-Based Access Control Model for Secure Big Data Processing in Hadoop Ecosystem,'' \textit{Proc. ACM Workshop on Attribute-Based Access Control}, pp. 13--24, 2020.

\bibitem{ye2014}
N. Ye, Y. Zhu, R. Wang, R. Malekian, and Q. Lin, ``An Efficient Authentication and Access Control Scheme for Perception Layer of Internet of Things,'' \textit{Appl. Math. Inf. Sci.}, vol. 8, no. 4, pp. 1617--1624, 2014.

\bibitem{patel2016}
V. M. Patel, R. Chellappa, D. Chandra, and B. Barbello, ``Continuous User Authentication on Mobile Devices: Recent Progress and Remaining Challenges,'' \textit{IEEE Signal Process. Mag.}, vol. 33, no. 4, pp. 49--61, Jul. 2016.

\bibitem{rasmussen2009}
K. B. Rasmussen and S. Capkun, ``Location Privacy of Distance Bounding Protocols,'' \textit{Proc. 15th ACM Conf. Computer and Communications Security (CCS)}, pp. 149--160, 2008.

\end{thebibliography}

\end{document}
